\scanpage{56}
% archive sheet number: 55

\handwritten{Supposons. d'abord $k$ algébriquement clos, $G$ réduit et connexe et $H$ réduit. D'après (\uline{BIBLE} \deleted{\unsure{VII|VIII}}\inserted{\deleted{Exp} 4, cor du th.\uline{1}\begin{tikzpicture}[x=1cm,y=1cm,baseline=-.1cm]\draw[thin] (0,-.1).. controls (-.1,-.35) and (-1.5,-.1)..(-1.7,-.4)--(-1.75,-.6)--(-1.8,-.35);\end{tikzpicture}}), il existe une representation lineaire de dimension finie $\rho:G\longrightarrow\mathrm{GL}(V)$ telle que $(\operatorname{Ker}\rho)_{\mathrm{red}}=H$. Par 11.11, $G/\operatorname{Ker}\rho$ est affine. donc $(G/H)/(\operatorname{Ker}\rho/H)$ est affine, donc aussi $G/H$, le groupe $\operatorname{Ker}\rho/H$ étant fini.}

\handwritten{Si $k$ est algébriquement clos, et $G$ et $H$ réduits, alors $(G/H)/(G^{\unsure{\circ|0|o}}/H\cap G^{\unsure{\circ|0|o}})$ est fini,\inserted{ comme quotient de $G/G^{\unsure{\circ|0|o}}$, et\begin{tikzpicture}[x=1cm,y=1cm,baseline=-.1cm]\draw[thin] (0,.05).. controls (0,-.25) and (-1.9,-.15)..(-2,-.35)--(-2.02,-.55);\end{tikzpicture}} $G^{\unsure{\circ|0|o}}/H\cap G^{\unsure{\circ|0|o}}$ est affine, donc $G/H$ est affine.}
\margin{\figure{\handwritten{\begin{tikzpicture}[x=0.2cm,y=0.4cm]\draw (0,8) .. controls (1.2,6.8) and (1,3) .. (1.2,0);\draw (0.8,6.8).. controls (0.7,6.5) and (0.3,6.5)..(-0.2,6.45);\end{tikzpicture}}}}

\figure{\handwritten{\begin{tikzpicture}\draw (0,0) .. controls (3,0.2) and (8,0.2) .. (12,0);\end{tikzpicture}}}
\deleted{\handwritten{\uline{N.B.}1) le th \unsure{entraîne|entraine} que si $G$ est affine algébrique et $H$ invariant, il existe une rep lin\unsure{e|é}aire de dim finie $\rho:G\longrightarrow\mathrm{GL}(V)$ telle que $H=\operatorname{Ker}\rho$}}

\deleted{\handwritten{2) La réf à BIBLE est un corollaire de 11.16, la d\unsure{e|é}monstration est facile mais un peu longue.}}

\deleted{\handwritten{3) Dan la réf à BIBLE Supposer $G$ connexe, car la dém. de BIBLE est \unsure{formulée} dans le cas general\begin{tikzpicture}[x=1cm,y=1cm,baseline=-.1cm]\draw[thin] (0,0).. controls (.45,.12) and (.85,.6)..(1.22,.5).. controls (1.43,.47) and (1.48,.25)..(1.4,.16);\end{tikzpicture}}}
\figure{\handwritten{\begin{tikzpicture}\draw (0,0) .. controls (4,0.4) and (8,0.5) .. (12,0.6);\end{tikzpicture}}}
